%%=============================================================================
%% Samenvatting
%%=============================================================================

% TODO: De "abstract" of samenvatting is een kernachtige (~ 1 blz. voor een
% thesis) synthese van het document.
%
% Een goede abstract biedt een kernachtig antwoord op volgende vragen:
%
% 1. Waarover gaat de bachelorproef?
% 2. Waarom heb je er over geschreven?
% 3. Hoe heb je het onderzoek uitgevoerd?
% 4. Wat waren de resultaten? Wat blijkt uit je onderzoek?
% 5. Wat betekenen je resultaten? Wat is de relevantie voor het werkveld?
%
% Daarom bestaat een abstract uit volgende componenten:
%
% - inleiding + kaderen thema
% - probleemstelling
% - (centrale) onderzoeksvraag
% - onderzoeksdoelstelling
% - methodologie
% - resultaten (beperk tot de belangrijkste, relevant voor de onderzoeksvraag)
% - conclusies, aanbevelingen, beperkingen
%
% LET OP! Een samenvatting is GEEN voorwoord!

%%---------- Nederlandse samenvatting -----------------------------------------
%
% TODO: Als je je bachelorproef in het Engels schrijft, moet je eerst een
% Nederlandse samenvatting invoegen. Haal daarvoor onderstaande code uit
% commentaar.
% Wie zijn bachelorproef in het Nederlands schrijft, kan dit negeren, de inhoud
% wordt niet in het document ingevoegd.

\IfLanguageName{english}{%
\selectlanguage{dutch}
\chapter*{Samenvatting}
\lipsum[1-4]
\selectlanguage{english}
}{}

%%---------- Samenvatting -----------------------------------------------------
% De samenvatting in de hoofdtaal van het document

\chapter*{\IfLanguageName{dutch}{Samenvatting}{Abstract}}

Veel grote organisaties vertrouwen vandaag nog op bedrijfskritische applicaties die ontwikkeld werden met
CA~Gen, een modelgedreven CASE-tool en vierde generatie ontwikkelomgeving uit de jaren tachtig. Hoewel
deze applicaties stabiel functioneren, neemt de beschikbare expertise zienderogen af en wordt onderhoud
steeds duurder en risicovoller. Binnen ArcelorMittal vormt de Batch Annealing Facility (BAF), een
deelmodule van het PLAPO-lijnplanningslandschap, de laatste CA~Gen-component die nog moet worden
gemigreerd naar PL/I. Eerdere migraties van zustermodules zijn weliswaar geslaagd, maar verspreid over
verschillende jaren en ontwikkelaars uitgevoerd, zonder dat een gestructureerde, herbruikbare aanpak werd
vastgelegd.

Dit onderzoek pakt die leemte aan door een \emph{high-level migratieframework} te ontwerpen voor de
gecontroleerde migratie van CA~Gen-applicaties naar PL/I. De centrale onderzoeksvraag luidt: \textit{``Hoe
kan een high-level framework worden opgesteld om CA~Gen-applicaties succesvol te migreren naar
PL/I-code?''} Het framework wordt opgebouwd rond vier kerncomponenten: (1)~een gestructureerde
analysefase voor het ontsluiten van CA~Gen-modellen, (2)~geconsolideerde vertaalrichtlijnen tussen
CA~Gen-constructen en PL/I-equivalenten, (3)~een set best practices en valkuilen, en (4)~een
validatiestrategie gebaseerd op de \emph{Parallel Run}-aanpak.

Het onderzoek is opgevat als een design-science traject waarbij het framework iteratief is opgebouwd op
basis van documentanalyse, literatuurstudie, reverse engineering van bestaande modules en een
proof-of-concept (PoC). Voor de PoC werd de CA~Gen-module \texttt{TPPBAMT} gemigreerd naar PL/I,
opgesplitst in een generieke MAIN-module, een specifieke BAF MAIN-module, een materiaalmodule, een
bestellingmodule en een output-generator. De resultaten werden gevalideerd door middel van een Parallel
Run waarbij de outputbestanden van het origineel en de migratie naast elkaar werden vergeleken.

De resultaten tonen aan dat het framework de migratie van \texttt{TPPBAMT} succesvol heeft begeleid en dat
de PoC functioneel equivalente output produceert. Het framework biedt ArcelorMittal een herhaalbare en
gedocumenteerde aanpak voor toekomstige CA~Gen-migraties en vermindert het risico op kennisverlies bij
ad-hoc trajecten. De analysefase, best-practices en validatiestrategie zijn taal-agnostisch en kunnen
hergebruikt worden voor migraties naar andere doeltalen, mits een aangepaste vertaaltabel.
